% GENERATED FILE. Edit canonical lessons, then run npm run generate:books.
% canonical-source-hash: fnv1a64:acbce47df22337de
% canonical-lessons: GE-C186-schneien, GE-C186-schuetzen, GE-C186-segeln, GE-C186-senden, GE-C186-spazieren

\chapter{To snow, To protect, To sail, To send, To stroll}
\label{ch:ge-to-snow-to-protect-to-sail-to-send-to-stroll}
\hlchaptermodality{186}

\begin{chapteropening}
\textbf{By the end of this chapter:} \emph{I can say to snow, to protect, to sail, to send and to stroll.}

Complete the last lesson of chapter 186: I can say to snow, to protect, to sail, to send and to stroll.
\end{chapteropening}

\section[schneien]{schneien — to snow}

\label{lesson:GE-C186-schneien}



\begin{quote}
Before the new one: say the German for to hit, then the German for to taste.
\end{quote}

\subsection*{You'll want to know: schneien}
\textbf{schneien} — \textquotedblleft{}to snow\textquotedblright{}.

\begin{cousinweb}
\textbf{schneien} $\leftarrow$ Old High German \emph{snīwan} $\leftarrow$ Proto-Germanic \emph{*snīwaną}, \textquotedblleft{}to snow\textquotedblright{} — inherited, from the same root as English \emph{snow}. (EtymDB 2.1)
\end{cousinweb}

\begin{grammarlens}[title={What you've built}]
One more word for this chapter.
\end{grammarlens}

\subsection*{Your turn}
\begin{itemize}
\raggedright
  \item \emph{Say it:} \emph{schneien}
  \item \emph{Say it:} \emph{schneien}, once more
  \item \emph{Say it:} say \emph{schneien}
  \item \emph{Recall:} say the German for to hit, then the German for to taste, then say \emph{schneien} again
\end{itemize}

\begin{tcolorbox}[breakable,colback=teal!4,colframe=teal!35!black,title={Before you move on}]
What does schneien mean? (\textquotedblleft{}To snow\textquotedblright{}.) Say it once more.
\end{tcolorbox}
\section[schützen]{schützen — to protect}

\label{lesson:GE-C186-schuetzen}



\begin{quote}
Before the new one: say the German for to taste, then the German for to snow.
\end{quote}

\subsection*{You'll want to know: schützen}
\textbf{schützen} — \textquotedblleft{}to protect\textquotedblright{}.

\begin{cousinweb}
\textbf{schützen} is the verb of the word for \textquotedblleft{}protection\textquotedblright{}, Middle High German \emph{schuz} (EtymDB 2.1): first \textquotedblleft{}to dam up water\textquotedblright{}, then \textquotedblleft{}to shield\textquotedblright{}.
\end{cousinweb}

\begin{grammarlens}[title={What you've built}]
One more word for this chapter.
\end{grammarlens}

\subsection*{Your turn}
\begin{itemize}
\raggedright
  \item \emph{Say it:} \emph{schützen}
  \item \emph{Say it:} \emph{schützen}, once more
  \item \emph{Say it:} say \emph{schützen}
  \item \emph{Recall:} say the German for to taste, then the German for to snow, then say \emph{schützen} again
\end{itemize}

\begin{tcolorbox}[breakable,colback=teal!4,colframe=teal!35!black,title={Before you move on}]
What does schützen mean? (\textquotedblleft{}To protect\textquotedblright{}.) Say it once more.
\end{tcolorbox}
\section[segeln]{segeln — to sail}

\label{lesson:GE-C186-segeln}



\begin{quote}
Before the new one: say the German for to snow, then the German for to protect.
\end{quote}

\subsection*{You'll want to know: segeln}
\textbf{segeln} — \textquotedblleft{}to sail\textquotedblright{}.

\begin{cousinweb}
\textbf{segeln} is the verb of the word for \textquotedblleft{}a sail\textquotedblright{}, Old High German \emph{segal} (EtymDB 2.1) $\leftarrow$ Proto-Germanic \emph{*seglą} — the same word as English \emph{sail}.
\end{cousinweb}

\begin{grammarlens}[title={What you've built}]
One more word for this chapter.
\end{grammarlens}

\subsection*{Your turn}
\begin{itemize}
\raggedright
  \item \emph{Say it:} \emph{segeln}
  \item \emph{Say it:} \emph{segeln}, once more
  \item \emph{Say it:} say \emph{segeln}
  \item \emph{Recall:} say the German for to snow, then the German for to protect, then say \emph{segeln} again
\end{itemize}

\begin{tcolorbox}[breakable,colback=teal!4,colframe=teal!35!black,title={Before you move on}]
What does segeln mean? (\textquotedblleft{}To sail\textquotedblright{}.) Say it once more.
\end{tcolorbox}
\section[senden]{senden — to send, to broadcast}

\label{lesson:GE-C186-senden}



\begin{quote}
Before the new one: say the German for to protect, then the German for to sail.
\end{quote}

\subsection*{You'll want to know: senden}
\textbf{senden} — \textquotedblleft{}to send, to broadcast\textquotedblright{}.

\begin{cousinweb}
\textbf{senden} $\leftarrow$ Middle High German \emph{senden} $\leftarrow$ Old High German \emph{senten} — inherited, the same verb as English \emph{send} and Gothic \emph{sandjan}. (EtymDB 2.1)
\end{cousinweb}

\begin{grammarlens}[title={What you've built}]
One more word for this chapter.
\end{grammarlens}

\subsection*{Your turn}
\begin{itemize}
\raggedright
  \item \emph{Say it:} \emph{senden}
  \item \emph{Say it:} \emph{senden}, once more
  \item \emph{Say it:} say \emph{senden}
  \item \emph{Recall:} say the German for to protect, then the German for to sail, then say \emph{senden} again
\end{itemize}

\begin{tcolorbox}[breakable,colback=teal!4,colframe=teal!35!black,title={Before you move on}]
What does senden mean? (\textquotedblleft{}To send, to broadcast\textquotedblright{}.) Say it once more.
\end{tcolorbox}
\section[spazieren]{spazieren — to stroll (spazieren gehen)}

\label{lesson:GE-C186-spazieren}



\begin{quote}
Before the new one: say the German for to sail, then the German for to send.
\end{quote}

\subsection*{You'll want to know: spazieren}
\textbf{spazieren} — \textquotedblleft{}to stroll (spazieren gehen)\textquotedblright{}.

\begin{cousinweb}
\textbf{spazieren} $\leftarrow$ Italian \emph{spaziare}, \textquotedblleft{}to roam\textquotedblright{} $\leftarrow$ Latin \emph{spatiārī}, \textquotedblleft{}to walk about\textquotedblright{}, from \emph{spatium}, \textquotedblleft{}space\textquotedblright{} (EtymDB 2.1) — borrowed.
\end{cousinweb}

\begin{grammarlens}[title={What you've built}]
One more word for this chapter.
\end{grammarlens}

\subsection*{Your turn}
\begin{itemize}
\raggedright
  \item \emph{Say it:} \emph{spazieren}
  \item \emph{Say it:} \emph{spazieren}, once more
  \item \emph{Say it:} say \emph{spazieren}
  \item \emph{Recall:} say the German for to sail, then the German for to send, then say \emph{spazieren} again
\end{itemize}

\begin{tcolorbox}[breakable,colback=teal!4,colframe=teal!35!black,title={Before you move on}]
What does spazieren mean? (\textquotedblleft{}To stroll (spazieren gehen)\textquotedblright{}.) Say it once more.
\end{tcolorbox}
